\section{Effective parameter spaces and the generic minimum}
\label{sec:effective-generic}

The auxiliary parameter spaces required in the preceding algorithm can
be constructed from equations. A second algorithm computes the generic
minimum when the target is given in an explicit Betti marking. The
regularity, Hilbert-scheme, and triangulation algorithms used here are
classical inputs; the conclusions below apply them to the particular
bounded-cycle minimum of this paper.

\begin{proposition}[R42: constructing an exhaustive bounded parameter space]
\label{prop:R42}
Let $X\hookrightarrow\mathbb P^n_S$ be a projective scheme given by
equations, where $S$ is the spectrum of an explicitly presented number
field, finite field, or localization of a number ring. Suppose the
embedding line bundle is $H^{\otimes m}$, for an ample integral line
bundle $H$ and a known positive integer $m$. Given $D\geq0$, one can
construct the projective schemes $Q_D$ and $H_D=Q_D\times_S Q_D$ of
B01, their universal families, and their degree labels, so that every
geometric effective integral surface cycle of $H$-degree at most $D$
occurs in $Q_D$. The construction is uniform in the residue
characteristic and includes the top-dimensional cycles of all boundary
fibers.
\end{proposition}

\begin{proof}
For $D=0$, take $Q_0=S$ with its zero cycle. Assume $D\geq1$, and put
\[
 E=m^2D,\qquad R=\left(\frac{E(E+1)}2\right)^2.
\]
If $Y\subset\mathbb P^n_{\bar k}$ is an integral surface of embedding
degree $e\leq E$, its homogeneous coordinate ring is reduced and
equidimensional of Krull dimension three, with arithmetic degree $e$.
Hoa's regularity theorem, valid over any field, gives
\[
 \operatorname{reg}I_Y
 \leq\left(\frac{e(e-1)}2+e\right)^{2^{3-2}}
 =\left(\frac{e(e+1)}2\right)^2\leq R
 \qquad\text{\cite[Theorem 0.2, author PDF p.~2]{HoaRegularity}}.
\]
The dimension in this theorem is the dimension of the homogeneous
coordinate ring, not the projective dimension of $Y$.

Let $P_Y$ be the Hilbert polynomial. Since
$\operatorname{reg}(\bar k[x_0,\ldots,x_n]/I_Y)\leq R-1$, the
Hilbert function agrees with $P_Y$ in every degree at least $R$.
Consequently its quadratic polynomial is determined by three integers
\[
 P_Y(R+j)=h_Y(R+j),\qquad
 0\leq h_Y(R+j)\leq\binom{n+R+j}{n},\qquad j=0,1,2.
\]
Enumerate every triple in these finite intervals and interpolate the
corresponding polynomial. Retain the quadratic polynomials whose
leading coefficient is $e/2$ for an integer $1\leq e\leq E$ such
that $e/m^2$ is a positive integer at most $D$. Every desired integral
support's Hilbert polynomial appears. Extraneous polynomials do no
harm: a polynomial which is not admissible gives an empty Hilbert
scheme, and every nonempty Hilbert scheme in the retained list still
has surface fibers of the recorded leading degree. Admissibility and
the Gotzmann construction are effective.

For each polynomial $P$ in this finite list, construct
$\operatorname{Hilb}^P(X/S)$ and its universal closed subscheme.
Over a field, this is precisely the algorithm of
Poonen--Testa--van Luijk, Lemma~8.23
\cite[author PDF p.~15]{PoonenTestaVanLuijk}. Over the indicated
arithmetic base, use the same integral Grassmannian equations:
Haiman--Sturmfels work over an arbitrary commutative ground ring, and
their Lemma~4.1, Proposition~4.2 and equations~(18)--(20) construct the
classical Hilbert scheme from the computable Gotzmann number
\cite[\S4, author PDF pp.~19--21]{HaimanSturmfels}.
Their Theorem~4.4 gives an alternative single-Grassmannian equation
system. Incidence with the explicitly given $X$ is imposed by requiring
its homogeneous ideal generators, multiplied to a sufficiently large
degree, to belong to the universal ideal subspace. These are finite
linear incidence equations in the Grassmannian charts. The universal
subscheme is obtained from that universal ideal. Thus both equations
and a finite affine cover can be computed, including over the number
ring.

Write $\delta(P)=e(P)/m^2$. Form the finite disjoint union of products
\[
 Q_D=\coprod_{\substack{0\leq q\leq D,\ 1\leq n_i\leq D\\
                        P_i\text{ in the finite list},\
                        \sum_i n_i\delta(P_i)\leq D}}
       \prod_{i=1}^q\operatorname{Hilb}^{P_i}(X/S),
\]
where the empty product supplies the zero cycle. Every effective cycle
$\sum_i n_i[Y_i]$ of degree at most $D$ occurs, because every
irreducible support has positive integral $H$-degree. Conversely,
every geometric point of this displayed scheme yields the weighted
sum of the top-dimensional fundamental cycles of its universal
fibers, with degree $\sum_i n_i\delta(P_i)\leq D$.
Its fibers may have embedded or lower-dimensional components; the
top-dimensional cycle and its degree remain valid, exactly as in B01.
Taking the product gives $H_D$. All lists, products, degree labels,
and universal families have therefore been constructed by finite
operations.
\end{proof}

\paragraph{The exact topological input.}
Fix a number field $K$ with a specified embedding into
$\overline{\mathbb Q}\subset\mathbb C$, described by algebraic root
data. Let $X/K$ be a smooth projective fourfold given by equations.
By an \emph{explicit Betti marking} we mean a finite simplicial
complex $K_0$, a semialgebraic homeomorphism
\[
                  h_0:|K_0|\longrightarrow X(\mathbb C)
\]
whose graph is given over the real algebraic numbers, and a rational
simplicial $4$-cocycle $\tau$ representing the target $t$ in this
marking. Equivalently, a basis of cohomology may be specified by
explicit cocycles on $K_0$ and $t$ by rational coordinates in that
basis. An abstract basis with no comparison map is insufficient input.
Semialgebraic triangulation provides such a topological model for
$X$; the target's coordinates in it are part of the input. Compare
\cite[Proposition 7.8, author PDF pp.~8--9]{PoonenTestaVanLuijk}.

\begin{theorem}[R43: computing the generic minimum from a marked class]
\label{thm:R43}
Let $X/K$ be an explicitly presented smooth projective fourfold,
with an explicit ample integral line bundle $H$ and an embedding
given by a known power of $H$. Supply an explicit Betti marking of a
rational class $t$ and an integer $C\geq1$ known to satisfy
$c_H(t)\leq C$. Then a terminating algorithm computes the exact
geometric minimum $C_0=c_H(t)$ and returns effective integral cycles
$Z^+,Z^-$ over a specified finite extension of $K$, together with
$N>0$, such that
\[
        [Z^+]-[Z^-]=Nt,\qquad
        \max(N,\deg_H Z^+,\deg_H Z^-)=C_0.
\]
The equality is in rational Betti cohomology. For an absolute Hodge
target with the compatible realizations of B02, this supplies the
actual global cycle needed in the simultaneous specialization
argument. In the Weil-type setting, R24 may provide $C$ when its
fixed lattice and polarization data are explicit.
\end{theorem}

\begin{proof}
Construct $Q_C$ by R42. Compute its finitely many geometric connected
components over $\overline K$, retaining equations over finite
extensions of $K$. This can be done by absolute decomposition and
the incidence graph of the resulting irreducible components:
nonemptiness of an intersection is an ideal-membership test on
projective affine charts. The same component operation is used in
\cite[Lemma 8.29(c)(iv), author PDF p.~17]{PoonenTestaVanLuijk}.
Choose an algebraic point on each nonempty connected component, by
solving its equations over $\overline K$. A finite collection of
such points and all their universal fibers are defined over a common
finite extension. From these fibers compute the top-dimensional
fundamental cycles, including the generic lengths giving any
multiplicities, and then apply the weights defining $Q_C$.
This yields finitely many explicit effective cycles $Z_j$, with
known degrees $d_j\leq C$.

By B01 the class of the universal cycle is constant on each
geometric connected component, and the degree label is constant
there as well. Every effective complex cycle of degree at most $C$
therefore has the same class and degree as some $Z_j$. In particular
this selection loses neither transcendental parameter points nor
boundary cycles. It also includes the zero cycle.

\emph{A common triangulation and comparison with the input marking.}
Represent projective complex space as the compact real semialgebraic
set of rank-one Hermitian projection matrices, using
$[z]\mapsto zz^*/(z^*z)$. This makes $X(\mathbb C)$, each reduced
support of each $Z_j$, its singular locus, and the images under $h_0$
of the simplices of $K_0$ explicit semialgebraic sets over the real
algebraic numbers. Quantifier elimination computes descriptions when
needed. The compatible triangulation algorithm gives a finite
complex $K$ and a semialgebraic homeomorphism
$h:|K|\to X(\mathbb C)$ adapted to all these subsets at once
\cite[Theorem 4.5, author PDF p.~14]{BasuTriangulation}.

The comparison with $K_0$ is effective, not merely an unspecified
isomorphism between abstract cohomology groups. For a vertex $v$ of
$K$, determine the unique open simplex $\sigma(v)$ of $K_0$
containing $h_0^{-1}h(v)$, and choose a vertex of its closure.
The resulting vertex map extends to a simplicial map
$F:K\to K_0$. Indeed every open simplex of $K$ maps under
$h_0^{-1}h$ into one open simplex of $K_0$, by compatibility, so the
images of all of its vertices belong to the closure of that simplex.
On each closed simplex, $|F|$ and $h_0^{-1}h$ take values in a common
closed simplex of $K_0$. The straight-line homotopy inside it proves
that $|F|$ is homotopic to $h_0^{-1}h$. All carrier decisions are
semialgebraic membership tests. Hence $F^*\tau$ represents exactly
the input class in the new triangulation.

\emph{Integral fundamental cycles.}
For each integral support $Y$ of a selected cycle, its inverse image
under $h$ is a subcomplex. Orient its top-dimensional simplices by
the complex orientation of $Y_{\rm sm}$, then add them with coefficient
one. The singular locus has real dimension at most two, so this
oriented four-chain is a cycle and is the integral fundamental class
of $Y$. Repeating with the computed multiplicities gives an integral
cycle chain for every $Z_j$.

The orientation choice is also an effective operation. In the interior
of a top-dimensional simplex choose a real algebraic point where the
semialgebraic parametrization is differentiable with full rank and
where an algebraic complex coordinate projection is an \emph{\'etale}
local chart on $Y_{\rm sm}$. Such points form a nonempty open subset:
a semialgebraic homeomorphism has a dense open differentiable locus
of full rank after stratification. The sign of its real Jacobian in
the ordered real and imaginary coordinates of this chart determines
the required orientation. The graph, differentiability and derivative
predicates are semialgebraic, so quantifier elimination and real
algebraic sign tests find such a point and determine that sign.
The same procedure orients the eight-dimensional fundamental cycle
$[X(\mathbb C)]$.

Simplicial boundary matrices and Smith normal form now compute the
integral homology classes $z_j\in H_4(K,\mathbb Z)$. The simplicial
cap product computes
\[
        t^\vee=(F^*\tau)\mathbin{\frown}[X(\mathbb C)]
                    \in H_4(K,\mathbb Q).
\]
Poincar\'e duality identifies $t^\vee$ with the target and each $z_j$
with its cycle class. Thus the exact rational equality
$z_i-z_j=Nt^\vee$ is decidable by finite linear algebra. No numerical
period comparison or Hodge-membership oracle is involved.

Finally, test this equality for every $1\leq N\leq C$ and every
ordered pair $i,j$, and take the smallest successful value of
$\max(N,d_i,d_j)$. The supplied upper bound guarantees a successful
test. Every test uses actual integral cycles; every possible
representation of complexity at most $C$ occurs among their component
classes. This proves both termination and exact minimality, and the
chosen universal fibers give equations for a representative attaining
the minimum.
\end{proof}

The input bound is used only to guarantee that this finite search
finds a representation. Without that hypothesis, the same bounded
procedure still decides whether a representation of complexity at
most $C$ exists. R43 does not assert that an arbitrary unnamed Betti
vector can be placed in the specified marking automatically, or that
the Hodge condition of an arbitrary rational vector is decidable.
The arithmetic continuation uses the effective integer of R44 and
the finite-field procedure R40; the parameter-space presentations are
provided by R42.
